% Options for packages loaded elsewhere
\PassOptionsToPackage{unicode}{hyperref}
\PassOptionsToPackage{hyphens}{url}
%
\documentclass[
]{article}
\usepackage{amsmath,amssymb}
\usepackage{iftex}
\ifPDFTeX
  \usepackage[T1]{fontenc}
  \usepackage[utf8]{inputenc}
  \usepackage{textcomp} % provide euro and other symbols
\else % if luatex or xetex
  \usepackage{unicode-math} % this also loads fontspec
  \defaultfontfeatures{Scale=MatchLowercase}
  \defaultfontfeatures[\rmfamily]{Ligatures=TeX,Scale=1}
\fi
\usepackage{lmodern}
\ifPDFTeX\else
  % xetex/luatex font selection
\fi
% Use upquote if available, for straight quotes in verbatim environments
\IfFileExists{upquote.sty}{\usepackage{upquote}}{}
\IfFileExists{microtype.sty}{% use microtype if available
  \usepackage[]{microtype}
  \UseMicrotypeSet[protrusion]{basicmath} % disable protrusion for tt fonts
}{}
\makeatletter
\@ifundefined{KOMAClassName}{% if non-KOMA class
  \IfFileExists{parskip.sty}{%
    \usepackage{parskip}
  }{% else
    \setlength{\parindent}{0pt}
    \setlength{\parskip}{6pt plus 2pt minus 1pt}}
}{% if KOMA class
  \KOMAoptions{parskip=half}}
\makeatother
\usepackage{xcolor}
\usepackage[margin=1in]{geometry}
\usepackage{graphicx}
\makeatletter
\newsavebox\pandoc@box
\newcommand*\pandocbounded[1]{% scales image to fit in text height/width
  \sbox\pandoc@box{#1}%
  \Gscale@div\@tempa{\textheight}{\dimexpr\ht\pandoc@box+\dp\pandoc@box\relax}%
  \Gscale@div\@tempb{\linewidth}{\wd\pandoc@box}%
  \ifdim\@tempb\p@<\@tempa\p@\let\@tempa\@tempb\fi% select the smaller of both
  \ifdim\@tempa\p@<\p@\scalebox{\@tempa}{\usebox\pandoc@box}%
  \else\usebox{\pandoc@box}%
  \fi%
}
% Set default figure placement to htbp
\def\fps@figure{htbp}
\makeatother
\setlength{\emergencystretch}{3em} % prevent overfull lines
\providecommand{\tightlist}{%
  \setlength{\itemsep}{0pt}\setlength{\parskip}{0pt}}
\setcounter{secnumdepth}{-\maxdimen} % remove section numbering
\usepackage{bookmark}
\IfFileExists{xurl.sty}{\usepackage{xurl}}{} % add URL line breaks if available
\urlstyle{same}
\hypersetup{
  pdftitle={README},
  hidelinks,
  pdfcreator={LaTeX via pandoc}}

\title{README}
\author{}
\date{\vspace{-2.5em}}

\begin{document}
\maketitle

\section{\texorpdfstring{\textbf{Rental Affordability Analysis: A
Comparison of Five European Capital Cities
(2022)}}{Rental Affordability Analysis: A Comparison of Five European Capital Cities (2022)}}\label{rental-affordability-analysis-a-comparison-of-five-european-capital-cities-2022}

\subsection{Project Goal}\label{project-goal}

This analysis investigates the rental affordability crisis across five
major European capital cities by calculating the Rent-to-Income Ratio
for an average single-person household. The goal is to determine which
city presents the most economically sensible option for a working
professional based on housing costs versus net monthly income.

\textbf{Key Insights \& Findings}

\textbf{The Outlier:} \emph{Lisbon} stands as the definitive
affordability outlier, consuming a shocking 76\% of the average monthly
net income for a 1-bedroom apartment.

\textbf{Most Favorable:} \emph{Madrid} and \emph{Berlin} showed the most
favorable affordability, requiring approximately 46-47\% of net income,
establishing them as the best options within the tested sample.

\emph{Income Paradox:} The analysis revealed that higher absolute income
(e.g., Berlin) does not guarantee proportionally higher affordability
compared to cities with lower incomes (e.g., Budapest), emphasizing the
criticality of the Rent-to-Income Ratio.

\emph{Visualization}: The findings are presented using a custom
Difference Stacked Bar Chart that clearly visualizes the rent burden on
top of the disposable income base.

\textbf{Data and Methodology}

\emph{Data Sources:} Eurostat (Household Income by NUTS 2 region) and
Eurostat (Current Market Rents Report). All data is sourced for 2022.

\emph{Metrics Calculated:} Monthly Net Income, Rent-to-Income Ratio, and
Disposable Income.

Tools: The entire analysis pipeline-from data cleaning and
transformation (using dplyr) to visualization (using ggplot2)-was
conducted in R.

\emph{Reproducibility:} The analysis is fully detailed in the
Rental\_Affordability\_Report.Rmd file.

\textbf{View the Full Report}

The complete, knitted HTML report, including all code, methodology, and
the final visualization, is available
\href{https://github.com/FatyaB/housing-affordability-analysis/blob/main/rent_affordability.html}{HERE}.

\textbf{Repository Contents}

Rental\_Affordability\_Report.Rmd: The complete R Markdown source file.

Rental\_Affordability\_Report.html: The final knitted report document
(Recommended viewing).

Housing Affordability Data Collection.xlsx: The raw input data used for
the analysis.

\end{document}
